\chapter{Interfaces That Remain Mathematical Research}\label{ch:frontier}

The selected reasoning chain is coherent where its maps share an explicit domain, operational semantics, and evidence language. That does not imply that every proposed component of a machine-native reasoning system has already been joined to every other one. It is useful to state what remains separate, because those boundaries suggest the most precise next theorems.

\section{Existence and acquisition are different problems}
The representation theorems identify when a smaller state \emph{can} preserve selected future observations. They do not generally tell an agent which smaller state to construct from sparse experience. Learning such a state requires a declared hypothesis family, a policy for obtaining observations, and conditions under which incompatible hypotheses can be separated. Statistical identifiability is not certain finite identification; every finite binary sequence can occur with positive probability under two distinct Bernoulli parameters in $(0,1)$.

A meaningful acquisition result would specify the loss or property of a proposed representation, the permitted query or intervention family, the distribution of observations and the tolerable failure probability. A failure certificate may show the available evidence cannot separate two task-relevant worlds. The resulting impossibility is relative to that evidence interface, not an unconditional impossibility of machine understanding.

\section{A learned semantic map is not a written commuting equation}
Conditional method transport has the form $TF_a=G_{\tau(a)}T$. This is meaningful when $T$, the operator translation $\tau$, and the applicable domains are already specified. A learned system must acquire them, ensure they cover its actual state trajectories, and show the diagrams continue to commute under the relevant interventions. A correct table for a finite declared task is a useful positive control, but it cannot certify that a large recurrent network learned a stable operator language across changed domains.

The Neuralese research documents both finite positive controls and failed stability or scientific gates. Those negative results delimit the next semantic acquisition contracts. An honest mathematical volume keeps the theorem about \emph{given} translators separate from an experiment asking whether one was discovered.

\section{A certificate is not a validated world action}
Proof-Path Theory defines typed justification fields, replayable traces and loss of information under compression. A checked proof records a consequence of its premises. Passing a certificate through an executable planning or action system requires additional semantics: which world facts correspond to the propositions, who may perform the selected action, and what actual state change the dispatcher implements.

The Typed Proof Learner and Resolution Completion branches retain relevant constructions, but the recovered handoff does not claim a complete connection from checked propositional derivation to world-action dispatch. Even a fully proved checker theorem would not automatically verify its external parser, message transport, or control interface. Those are additional refinement obligations.

\section{Revision requires dependency completeness}
A revision may change the observation vocabulary, the task, the permitted actions, the model family or the representation used to express an old result. The statement of a theorem can remain true while its application becomes invalid. A retained dependency graph makes selective invalidation possible if the graph is complete for the applicable premises. A hash proves byte identity, not the adequacy of the recorded dependency set.

The next composition theorem should therefore carry explicit compatibility data between an old and new contract, a translation of certificate premises, and a proof that each retained conclusion is still justified. In the absence of those data, an old result should remain accessible as historical mathematics but should not automatically authorize a new action. Formalizing that rule is more valuable than assigning another name to a ``self-correcting agent.''

\section{Resources, workspaces, and authority}
Improvement claims also require a common cost ledger. Reduced operator calls may be outweighed by the preparation, validation or storage needed to obtain a reusable method. Different price weights can change the sign of a proposed saving even when semantic correctness remains unchanged. The reviewed core already formalizes a finite version of this principle. Broad claims of scalable or indefinite improvement require additional assumptions about task distributions, resource limits and repeated acquisition.

Likewise, jointly reconstructible information is not sufficient for a functioning workspace or a machine society. Causal relevance, persistence, identity, access rules, permissions, institutional formation and accountable endings require their own theorems. The safe-ending calculation of Chapter~\ref{ch:endinganalysis} uses a stipulated action system. It does not show that authority was legitimately established or that real institutions will comply with a command.

\section{A bounded agenda}
The strongest next steps are mathematical and testable. Complete the four selected Lean correspondence targets. Prove further typed compositions between stochastic correction and evidence-preserving revision. Add explicit no-go examples for proposed learned semantic translators when an observation interface is insufficient. Then establish restricted refinement theorems linking a certificate checker to a producer and to a declared runtime, rather than promising end-to-end verification of every historical component. These are substantial research tasks, but each admits finite exact examples, sharp hypotheses and ordinary mathematical review before any large-scale training campaign.
